% !TEX program = pdflatex
\documentclass[11pt]{article}

% ── Layout ────────────────────────────────────────────────────────────────────
\usepackage[margin=1in]{geometry}
\usepackage{microtype}
\usepackage{parskip}
\usepackage{float}

% ── Fonts ─────────────────────────────────────────────────────────────────────
\usepackage[T1]{fontenc}
\usepackage{lmodern}
\usepackage{newtxtext}
\usepackage[scaled=0.92]{inconsolata}

% ── Math ──────────────────────────────────────────────────────────────────────
\usepackage{amsmath}
\usepackage{amssymb}

% ── Graphics, tables ──────────────────────────────────────────────────────────
\usepackage{graphicx}
\usepackage{xcolor}
\usepackage{booktabs}
\usepackage{array}
\usepackage{caption}
\usepackage{colortbl}

% ── Boxes ─────────────────────────────────────────────────────────────────────
\usepackage[most]{tcolorbox}

% ── Hyperlinks ────────────────────────────────────────────────────────────────
\usepackage{hyperref}

% ── Colors ────────────────────────────────────────────────────────────────────
\definecolor{accent}{HTML}{2E5A88}
\definecolor{accentred}{HTML}{C44E52}
\definecolor{softgray}{HTML}{F5F5F7}
\definecolor{rule}{HTML}{B7C2D0}
\definecolor{codebg}{HTML}{F2F4F7}
\definecolor{cellhi}{HTML}{E8F0F8}
\definecolor{cellwarn}{HTML}{F8E8E8}

% ── Section formatting ────────────────────────────────────────────────────────
\usepackage{titlesec}
\titleformat{\section}
  {\color{accent}\Large\bfseries\sffamily}
  {\thesection}{0.6em}{}
  [\vspace{-2pt}{\color{rule}\titlerule[0.6pt]}]
\titleformat{\subsection}
  {\color{accent}\large\bfseries\sffamily}
  {\thesubsection}{0.5em}{}
\titlespacing*{\section}{0pt}{1.4em}{0.6em}

% ── Captions ──────────────────────────────────────────────────────────────────
\captionsetup{font=small, labelfont={bf,color=accent}, skip=4pt}

% ── Hyperref ──────────────────────────────────────────────────────────────────
\hypersetup{colorlinks=true, linkcolor=accent, urlcolor=accent, citecolor=accent}

% ── Inline code ───────────────────────────────────────────────────────────────
\newcommand{\code}[1]{%
  \tcbox[on line, colback=codebg, colframe=codebg,
         boxrule=0pt, arc=2pt, left=2pt, right=2pt, top=0pt, bottom=0pt]
        {\ttfamily\small #1}}

% ── Callout box ───────────────────────────────────────────────────────────────
\newtcolorbox{keybox}[1][]{
  enhanced, breakable,
  colback=softgray, colframe=accent,
  boxrule=0pt, leftrule=3pt,
  arc=2pt, left=10pt, right=10pt, top=6pt, bottom=6pt,
  fonttitle=\bfseries\sffamily\color{accent},
  #1
}

\title{\vspace{-2.0em}{\color{accent}\textsf{\bfseries The Rand Index and the Adjusted Rand Index, from the ground up}}}
\author{}
\date{}

\begin{document}
\maketitle
\vspace{-3.0em}

\noindent
\textit{This note builds RI and ARI from first principles. No leaps, no
``it can be shown''. Every quantity is computed on small examples that you
can verify by hand. By the end you will see that ARI is not a mysterious
formula --- it is the unique answer to one very specific question, and you
will know which question.}

% =============================================================================
\section{The setup: comparing two groupings of the same items}

We have $N$ items (in our project: articles). Two different ``graders'' have
each sorted them into clusters.

\begin{itemize}\setlength\itemsep{1pt}
  \item The \textbf{truth} $T$: the real authorship labels.
        Truth partitions the items into clusters
        $T_1, T_2, \ldots, T_K$ of sizes $\alpha_1, \alpha_2, \ldots, \alpha_K$.
  \item The \textbf{prediction} $P$: what K-means produced.
        Prediction partitions the items into clusters
        $P_1, P_2, \ldots, P_L$ of sizes $\beta_1, \beta_2, \ldots, \beta_L$.
\end{itemize}

We want a single number that measures how well $P$ matches $T$. Higher
should mean better. That is the only goal of this entire document.

\paragraph{The first thing that goes wrong.}
The obvious idea --- ``count how many items have the same label in $T$ and
$P$'' --- does not work, because the cluster labels are arbitrary names.
If truth calls a cluster ``A'' and the prediction calls the same cluster
``7'', the prediction is perfect even though no item shares a name with the
truth. And it gets worse: K-means might split true cluster A into two
predicted clusters, or merge two true clusters into one --- there is no
clean re-labelling that fixes this.

% =============================================================================
\section{The trick: don't compare labels. Compare relationships.}

Here is the move that makes everything work.

Forget about cluster names entirely. For every \emph{pair} of items
$(i, j)$, $T$ makes one of two statements:
\begin{itemize}\setlength\itemsep{0pt}
  \item ``$i$ and $j$ are in the \textbf{same} true cluster'', or
  \item ``$i$ and $j$ are in \textbf{different} true clusters''.
\end{itemize}
$P$ makes one of the same two statements. These statements do not depend
on what the clusters are called. They only depend on the \emph{partition},
i.e.\ the grouping itself.

So we ask one question, repeated over every pair of items:
\begin{quote}\itshape
  Do $T$ and $P$ make the same statement about this pair?
\end{quote}

If we just count how often the answer is yes, we get a similarity score
that is immune to label-renaming. That count, divided by the total number
of pairs, is the Rand Index.

% =============================================================================
\section{The four kinds of pairs}

For each pair $(i, j)$ there are four possibilities, depending on what $T$
says and what $P$ says. We give each combination a letter --- this is the
notation everyone uses.

\begin{table}[H]
\centering
\renewcommand{\arraystretch}{1.5}
\begin{tabular}{|l|l|l|}
\hline
 & \textbf{Pred says SAME} & \textbf{Pred says DIFF} \\ \hline
\textbf{Truth says SAME}
  & \cellcolor{cellhi}$a$: both say together (true positive)
  & \cellcolor{cellwarn}$b$: pred wrongly split a true cluster \\ \hline
\textbf{Truth says DIFF}
  & \cellcolor{cellwarn}$c$: pred wrongly merged two true clusters
  & \cellcolor{cellhi}$d$: both say apart (true negative) \\ \hline
\end{tabular}
\end{table}

The shaded cells $a$ and $d$ are the \textbf{agreements}; the unshaded
cells $b$ and $c$ are the \textbf{disagreements}. Every pair of items
lands in exactly one of the four cells, so
\[
  a + b + c + d \;=\; \binom{N}{2}.
\]

% =============================================================================
\section{A small worked example}
\label{sec:example}

Let us pick a tiny example and stick with it for the rest of the document.

Take $N = 6$ items labelled $1, 2, 3, 4, 5, 6$. Truth and prediction:
\[
  T \;=\; \bigl\{\{1, 2, 3\},\; \{4, 5, 6\}\bigr\}, \qquad
  P \;=\; \bigl\{\{1, 2, 4\},\; \{3, 5, 6\}\bigr\}.
\]
So truth groups the first three together and the last three together;
the prediction is the same except items $3$ and $4$ are swapped.

Total pairs: $\binom{6}{2} = 15$. Let us list every single one.

\begin{table}[H]
\centering
\small
\begin{tabular}{cccc}
\toprule
Pair & In same true cluster? & In same pred cluster? & Category \\
\midrule
$(1, 2)$ & yes ($T_1$) & yes ($P_1$) & $a$ \\
$(1, 3)$ & yes ($T_1$) & no & $b$ \\
$(1, 4)$ & no & yes ($P_1$) & $c$ \\
$(1, 5)$ & no & no & $d$ \\
$(1, 6)$ & no & no & $d$ \\
$(2, 3)$ & yes ($T_1$) & no & $b$ \\
$(2, 4)$ & no & yes ($P_1$) & $c$ \\
$(2, 5)$ & no & no & $d$ \\
$(2, 6)$ & no & no & $d$ \\
$(3, 4)$ & no & no & $d$ \\
$(3, 5)$ & no & yes ($P_2$) & $c$ \\
$(3, 6)$ & no & yes ($P_2$) & $c$ \\
$(4, 5)$ & yes ($T_2$) & no & $b$ \\
$(4, 6)$ & yes ($T_2$) & no & $b$ \\
$(5, 6)$ & yes ($T_2$) & yes ($P_2$) & $a$ \\
\bottomrule
\end{tabular}
\end{table}

Tallying the column on the right:
\[
  a = 2, \qquad b = 4, \qquad c = 4, \qquad d = 5,
  \qquad a + b + c + d = 15 = \binom{6}{2}. \;\checkmark
\]

% =============================================================================
\section{The Rand Index}

\begin{keybox}
\[
  \text{RI}(T, P) \;=\; \frac{a + d}{\,a + b + c + d\,} \;=\; \frac{a + d}{\binom{N}{2}}.
\]
\end{keybox}

In words: the fraction of pairs on which $T$ and $P$ make the same
statement. On our example,
\[
  \text{RI} \;=\; \frac{2 + 5}{15} \;=\; \frac{7}{15} \;\approx\; 0.467.
\]

Range: $[0, 1]$. RI $= 1$ if and only if the two partitions are identical
(every pair agrees); RI $= 0$ if every pair disagrees, which is rare and
hard to construct.

So far, so good. The Rand Index is a clean, label-free way of measuring
how close two partitions are. The problem is what value it gives in
practice. We turn to that now.

% =============================================================================
\section{Problem 1: random clusterings score embarrassingly high}
\label{sec:problem1}

Here is the first thing that should bother you: an utterly useless
clustering does \emph{not} score RI $= 0$. It scores something like $0.5$,
and the exact number depends on how the problem is set up.

\paragraph{A concrete experiment we can do by hand.}
Take a clean toy problem: $N = 30$ items, truth has two clusters of $15$
each. Now I propose the dumbest possible prediction: independently for
each item, flip a fair coin --- heads goes into predicted cluster $1$,
tails into predicted cluster $2$. The coin doesn't know anything about
authorship. The prediction is pure noise.

What RI do we expect on average?

For any specific pair of items $(i, j)$, the prediction puts them in the
same cluster exactly when both coin flips agree (HH or TT), which happens
with probability $\tfrac{1}{2}$. This is independent of whether truth puts
them together --- the coins do not consult the truth.

\paragraph{Counting the pairs.}
\begin{itemize}\setlength\itemsep{1pt}
  \item \textbf{True-same pairs}: items within the same true cluster.
        Each true cluster of $15$ contributes $\binom{15}{2} = 105$ pairs,
        and there are two clusters, so $2 \cdot 105 = 210$ true-same pairs.
  \item \textbf{True-diff pairs}: every other pair.
        $\binom{30}{2} - 210 = 435 - 210 = 225$.
\end{itemize}

\paragraph{Expected agreements.}
For each true-same pair, the coin agrees (says ``same'') with probability
$\tfrac{1}{2}$. So:
\[
  E[a] \;=\; 210 \cdot \tfrac{1}{2} \;=\; 105.
\]
For each true-diff pair, the coin agrees (says ``diff'') with probability
$\tfrac{1}{2}$. So:
\[
  E[d] \;=\; 225 \cdot \tfrac{1}{2} \;=\; 112.5.
\]
And therefore:
\[
  E[\text{RI}] \;=\; \frac{E[a] + E[d]}{\binom{30}{2}}
  \;=\; \frac{105 + 112.5}{435} \;=\; 0.500.
\]

\medskip
\noindent
\textbf{So a coin-flip clustering scores RI $\approx 0.5$.} It learned
nothing, but half the pairs agreed with truth by accident. This means
RI = 0.5 is \emph{not} the midpoint between worst and best. It is the
baseline. You would need RI well above $0.5$ before claiming you have
learned anything from the data.

This is annoying but not yet fatal. We will see the killer problem next.

% =============================================================================
\section{Problem 2: the random baseline depends on $K$}
\label{sec:problem2}

Here is the deal-breaker. The number $0.5$ we just computed is the
baseline only for $K = 2$. It depends sensitively on how many clusters are
in the truth. As $K$ grows, the baseline grows toward $1$, and RI loses
its ability to discriminate.

\paragraph{The general setup.}
Suppose truth has $K$ clusters of equal size $M$ each, so $N = K M$. The
prediction is again pure noise: each item dropped uniformly at random into
one of $K$ predicted clusters. For any pair of items,
\[
  P(\text{pred says same}) \;=\; \frac{1}{K},
  \qquad
  P(\text{pred says diff}) \;=\; \frac{K - 1}{K}.
\]
(Because both items must land in the same one of $K$ predicted clusters,
which happens with probability $1/K$.)

\paragraph{Counting the true-same and true-diff pairs.}
True-same pairs: $K \cdot \binom{M}{2}$.
True-diff pairs: $\binom{N}{2} - K \binom{M}{2}$.

\paragraph{Expected agreements.}
\[
  E[a] \;=\; K \binom{M}{2} \cdot \frac{1}{K} \;=\; \binom{M}{2},
\]
\[
  E[d] \;=\; \Bigl[\binom{N}{2} - K \binom{M}{2}\Bigr] \cdot \frac{K - 1}{K},
\]
\[
  E[\text{RI}] \;=\; \frac{E[a] + E[d]}{\binom{N}{2}}.
\]

Now I plug in our project's setting --- $M = 15$ articles per author ---
and sweep $K$. Each row of the table below is the result of one full
calculation using the formulas above. Try one row yourself to verify; the
arithmetic takes 30 seconds.

\begin{table}[H]
\centering
\small
\begin{tabular}{r r r r r r r}
\toprule
$K$ & $N = KM$ & $\binom{N}{2}$ & true-same & true-diff & $E[a]$ & $E[d]$ \\
\midrule
2  & 30  & 435    & 210  & 225    & 105 & 112.5 \\
4  & 60  & 1{,}770   & 420  & 1{,}350  & 105 & 1{,}012.5 \\
6  & 90  & 4{,}005   & 630  & 3{,}375  & 105 & 2{,}812.5 \\
8  & 120 & 7{,}140   & 840  & 6{,}300  & 105 & 5{,}512.5 \\
10 & 150 & 11{,}175  & 1{,}050 & 10{,}125 & 105 & 9{,}112.5 \\
12 & 180 & 16{,}110  & 1{,}260 & 14{,}850 & 105 & 13{,}612.5 \\
\bottomrule
\end{tabular}
\caption{Components of $E[\text{RI}]$ under a random $K$-cluster prediction,
with $M = 15$ items per true cluster. Note that $E[a]$ stays constant at
$\binom{M}{2} = 105$ while $E[d]$ explodes.}
\end{table}

\noindent
Dividing $E[a] + E[d]$ by $\binom{N}{2}$ gives the random baseline:

\begin{table}[H]
\centering
\small
\begin{tabular}{r c l}
\toprule
$K$ & $E[\text{RI}]$ for random & what RI $\approx 0.85$ would mean \\
\midrule
2  & 0.500 & far above chance: clustering is genuinely working \\
4  & 0.631 & well above chance \\
6  & 0.729 & clearly above chance \\
8  & 0.787 & somewhat above chance \\
10 & 0.825 & barely above chance \\
12 & 0.852 & \emph{at} chance: a fair coin would do this \\
\bottomrule
\end{tabular}
\caption{The same RI score means wildly different things at different $K$.}
\label{tab:rkillertable}
\end{table}

\begin{keybox}[title={The killer}]
A single RI value, like $0.85$, can simultaneously mean ``great
performance'' (at $K = 2$) and ``literally random'' (at $K = 12$). If you
plot RI vs $K$ for a clustering algorithm of \emph{constant skill}, RI
will trend \emph{upward} purely because the random baseline trends
upward. The actual quality could be flat, climbing, or collapsing, and
you cannot tell which from RI alone.
\end{keybox}

\paragraph{Why does the baseline climb with $K$?}
Look at the table. At $K = 12$ there are $14{,}850$ true-diff pairs out of
$16{,}110$ total --- that's $92\%$. A random prediction also produces
many small predicted clusters, so the vast majority of its pairs are
pred-diff too. The two graders agree on ``these two items are different''
\emph{trivially}, because almost everything is different from almost
everything. That flood of trivial $d$-agreements inflates RI.

At $K = 2$, by contrast, ``true-diff'' is only about half of pairs, so
random doesn't get easy points for free.

The lesson: \textbf{RI rewards a clustering not for what it gets right,
but for how lopsided the partition is}. Lopsided structures
automatically have high RI baselines.

% =============================================================================
\section{The fix, conceptually}

We want to throw away the part of RI that comes from chance. Conceptually,
once we know the random baseline $E[\text{RI}]$:

\begin{itemize}\setlength\itemsep{1pt}
  \item Anything below $E[\text{RI}]$ is worse than random and should
        score negative.
  \item Exactly $E[\text{RI}]$ should score zero.
  \item RI $= 1$ should score one.
\end{itemize}

The simplest way to achieve this is a linear rescaling: subtract the
baseline, then divide by the distance from the baseline to the maximum.

\begin{keybox}[title={The conceptual ARI}]
\[
  \text{ARI} \;=\; \frac{\text{RI} - E[\text{RI}]}{1 - E[\text{RI}]}
  \;=\; \frac{\text{how much better than chance}}{\text{how much better
  perfect would be than chance}}.
\]
\end{keybox}

In words: \emph{ARI is the fraction of the gap from random to perfect
that your clustering closed.}

That sentence is the whole conceptual picture. The actual H--A formula
people write down differs slightly in technical details, but it answers
exactly this question. We need two more ingredients before we can write
it cleanly:

\begin{enumerate}\setlength\itemsep{1pt}
  \item A precise definition of ``random'', so we can compute
        $E[\text{RI}]$.
  \item A small algebraic rearrangement that makes the formula prettier
        and more numerically stable.
\end{enumerate}

% =============================================================================
\section{What does ``random'' mean? The Hubert--Arabie null}
\label{sec:null}

In Section~\ref{sec:problem1} we used a very informal random model: ``flip
a coin for each item''. That worked for the calculation but leaves a
fundamental choice unspecified: does our random predictor get to know the
true cluster sizes? Does it get to pick its own cluster sizes? What if
those choices accidentally match the truth's cluster sizes?

The standard answer, and the one ARI uses, is the
\textbf{Hubert--Arabie null model} (1985). It is the natural choice and we
will see why.

\paragraph{The null model.}
Take both partitions' cluster sizes as fixed.
\begin{itemize}\setlength\itemsep{1pt}
  \item Truth has clusters of sizes $\alpha_1, \alpha_2, \ldots, \alpha_K$
        --- these are observed and we hold them constant.
  \item Pred has clusters of sizes $\beta_1, \beta_2, \ldots, \beta_L$
        --- also observed and held constant.
\end{itemize}
Under the null, the only thing that is random is \emph{which items get
which combination of (true, pred) labels}, subject to the row and
column-size constraints.

Equivalently: build a contingency table whose row sums are
$(\alpha_i)$ and column sums are $(\beta_j)$, and fill in the cells
uniformly at random subject to those marginals.

\paragraph{Why this null?}
Three reasons.
\begin{enumerate}\setlength\itemsep{1pt}
  \item \textbf{It doesn't give you free credit for matching cluster
        sizes by accident.} The marginals are pinned to what we observed.
        Only the cross-tabulation is random.
  \item \textbf{It is parameter-free.} No tuning constants. Both
        $\alpha_i$ and $\beta_j$ are data.
  \item \textbf{It is a classical model in statistics} --- the same null
        behind Fisher's exact test. The math is well-understood.
\end{enumerate}

% =============================================================================
\section{The expected $a$ under the null}

To compute the expectation, we need to think about $a$ in terms of the
contingency table.

\paragraph{Re-deriving $a$ from the table.}
Let $n_{ij}$ be the count in cell $(i, j)$ --- the number of items that
are in true cluster $i$ and predicted cluster $j$. Two items
\emph{both} land in the same true cluster and the same predicted cluster
exactly when they sit in the same cell. The number of pairs inside cell
$(i,j)$ is $\binom{n_{ij}}{2}$, so:
\[
  a \;=\; \sum_{i, j} \binom{n_{ij}}{2}.
\]
This is just a rewrite of the definition of $a$, but it's the rewrite we
need: now $a$ is a sum over cells, and each cell's contribution depends
only on a count $n_{ij}$ whose distribution we know.

\paragraph{The distribution of $n_{ij}$ under the null.}
Under the Hubert--Arabie null, the row marginals are $\alpha_i$ and the
column marginals are $\beta_j$. The marginal distribution of $n_{ij}$ is
hypergeometric, with mean
$\,E[n_{ij}] \;=\; \alpha_i \, \beta_j / N\,$
(natural: row $i$ has $\alpha_i$ items, each independently sitting in
column $j$ with probability $\beta_j / N$).

A standard identity for the hypergeometric distribution gives the second
moment we actually need:
\[
  E\!\left[\binom{n_{ij}}{2}\right]
  \;=\; \frac{\binom{\alpha_i}{2}\binom{\beta_j}{2}}{\binom{N}{2}}.
\]
You can prove this by writing $\binom{n_{ij}}{2}$ as the number of pairs
of items both landing in cell $(i, j)$, and computing the probability for
each pair separately --- I will skip the proof here, but the formula is
classical and you can take it on faith.

\paragraph{Summing over all cells.}
\begin{align*}
  E[a] \;&=\; \sum_{i, j} E\!\left[\binom{n_{ij}}{2}\right]
       \;=\; \sum_{i, j} \frac{\binom{\alpha_i}{2}\binom{\beta_j}{2}}{\binom{N}{2}} \\
       \;&=\; \frac{1}{\binom{N}{2}} \left(\sum_i \binom{\alpha_i}{2}\right)
                                     \left(\sum_j \binom{\beta_j}{2}\right)
\end{align*}
where in the last step we used $\sum_{ij} x_i y_j = (\sum_i x_i)(\sum_j y_j)$.

\paragraph{Two clean quantities to name.}
Each of those sums has a meaning of its own.

\begin{keybox}
\[
  S_T \;=\; \sum_{i} \binom{\alpha_i}{2}
  \;=\; \text{(number of true-same pairs)}.
\]
\[
  S_P \;=\; \sum_{j} \binom{\beta_j}{2}
  \;=\; \text{(number of pred-same pairs)}.
\]
Then
\[
  E[a] \;=\; \frac{S_T \cdot S_P}{\binom{N}{2}}.
\]
\end{keybox}

That's it. $S_T$ is a property of truth alone; $S_P$ is a property of
prediction alone; their product divided by $\binom{N}{2}$ is the
chance-level number of TS--PS pairs we'd expect if there were no
relationship between them.

% =============================================================================
\section{The actual ARI formula}

Recall the conceptual ARI:
\[
  \text{ARI} \;=\; \frac{\text{observed agreement above chance}}
                       {\text{maximum possible agreement above chance}}.
\]
Hubert and Arabie wrote this using $a$ directly rather than RI, which
gives an equivalent but cleaner formula:

\begin{keybox}[title={Hubert--Arabie ARI}]
\[
  \text{ARI} \;=\; \frac{a - E[a]}{\mathrm{MaxIndex} - E[a]},
\]
where
\[
  E[a] = \frac{S_T \cdot S_P}{\binom{N}{2}}, \qquad
  \mathrm{MaxIndex} = \frac{S_T + S_P}{2}.
\]
\end{keybox}

\paragraph{Why is using $a$ equivalent to using RI?}
A useful identity: among the four counts $(a, b, c, d)$, only one is
independent once the marginals are fixed. Specifically, $a + b = S_T$
(every true-same pair is either TS--PS or TS--PD), so $b = S_T - a$.
Similarly $a + c = S_P$, so $c = S_P - a$. And $d = \binom{N}{2} - a - b - c$.
This means $b, c, d$ are all linear functions of $a$ (once marginals are
fixed), and so is RI. Adjusting $a$ for chance and adjusting RI for chance
give equivalent formulas, differing only by an overall scale that cancels
in the ratio.

\paragraph{Why MaxIndex $= (S_T + S_P)/2$?}
The largest $a$ can ever be is $\min(S_T, S_P)$ --- you can't have more
TS--PS pairs than the total of either type. So the
\emph{exact} maximum is $\min(S_T, S_P)$, and you could in principle put
that in the denominator. H--A use the average $(S_T + S_P)/2$ instead.

Why? It is for algebraic convenience --- the formula derived from
$E[\binom{n_{ij}}{2}]$ matches the H--A denominator exactly, so the whole
formula can be written cleanly in terms of $S_T$ and $S_P$. With the
average:
\begin{itemize}\setlength\itemsep{1pt}
  \item ARI $= 1$ iff $T$ and $P$ are identical partitions (so
        $S_T = S_P = a$, and both numerator and denominator equal
        $S_T - E[a]$).
  \item ARI can be slightly negative for worse-than-random clusterings,
        which is informative and not a bug.
  \item ARI is symmetric in $T$ and $P$, as it should be.
\end{itemize}

% =============================================================================
\section{Putting it all together: ARI for our running example}

Back to the $N = 6$ example from Section~\ref{sec:example}:
$T = \{\{1,2,3\}, \{4,5,6\}\}$,
$P = \{\{1,2,4\}, \{3,5,6\}\}$.

\paragraph{Step 1: build the contingency table.}
For each item we look at its $(T, P)$ pair of labels and tally.

\begin{table}[H]
\centering
\small
\begin{tabular}{c | c c | c}
& $P_1 = \{1, 2, 4\}$ & $P_2 = \{3, 5, 6\}$ & row total $\alpha_i$ \\
\hline
$T_1 = \{1, 2, 3\}$ & $2$ \; (items 1, 2) & $1$ \; (item 3) & $3$ \\
$T_2 = \{4, 5, 6\}$ & $1$ \; (item 4) & $2$ \; (items 5, 6) & $3$ \\
\hline
col total $\beta_j$ & $3$ & $3$ & $6 = N$ \\
\end{tabular}
\end{table}

\paragraph{Step 2: read off $a$ from the table.}
\[
  a \;=\; \sum_{i,j} \binom{n_{ij}}{2}
       \;=\; \binom{2}{2} + \binom{1}{2} + \binom{1}{2} + \binom{2}{2}
       \;=\; 1 + 0 + 0 + 1 \;=\; 2.
\]
(Matches what we counted by hand in Section~\ref{sec:example}.)

\paragraph{Step 3: compute $S_T$ and $S_P$.}
\[
  S_T \;=\; \binom{3}{2} + \binom{3}{2} \;=\; 3 + 3 \;=\; 6.
\]
\[
  S_P \;=\; \binom{3}{2} + \binom{3}{2} \;=\; 3 + 3 \;=\; 6.
\]

\paragraph{Step 4: compute $E[a]$.}
\[
  E[a] \;=\; \frac{S_T \cdot S_P}{\binom{N}{2}}
        \;=\; \frac{6 \cdot 6}{15} \;=\; \frac{36}{15} \;=\; 2.4.
\]
Interpretation: with these marginals (two clusters of three on each side),
a random arrangement would, on average, produce $2.4$ TS--PS pairs. We
observed $2$. So we're slightly \emph{below} expectation.

\paragraph{Step 5: compute MaxIndex and ARI.}
\[
  \mathrm{MaxIndex} \;=\; \frac{S_T + S_P}{2} \;=\; \frac{6 + 6}{2} \;=\; 6.
\]
\[
  \text{ARI} \;=\; \frac{a - E[a]}{\mathrm{MaxIndex} - E[a]}
         \;=\; \frac{2 - 2.4}{6 - 2.4} \;=\; \frac{-0.4}{3.6}
         \;\approx\; -0.111.
\]

\paragraph{Sanity-checking the negative value.}
RI for this example was $0.467$, which sounds OK. But ARI says we are
\emph{worse} than chance. Both can be true: under the H--A null with
these specific marginals (two clusters of three on each side), the
average random arrangement actually produces $2.4$ TS--PS pairs --- more
than the $2$ we observed. So our particular swap (item 3 and item 4) does
worse than typical noise.

This is the value of ARI: it tells you not just ``how many pairs match''
but ``how many pairs match \emph{compared to what nothing-going-on would
have produced}''.

% =============================================================================
\section{Three more sanity checks}

To make sure the formula is doing what we want, let's run it on three
quick examples (still $N = 6$, same truth $T$ as before).

\subsection{Perfect prediction}
$P = T = \{\{1,2,3\}, \{4,5,6\}\}$. Contingency table:
$\bigl[\begin{smallmatrix} 3 & 0 \\ 0 & 3 \end{smallmatrix}\bigr]$.
\[
  a = \binom{3}{2} + \binom{3}{2} = 6, \quad
  S_T = 6, \quad S_P = 6, \quad E[a] = 2.4, \quad \mathrm{MaxIndex} = 6.
\]
\[
  \text{ARI} = \frac{6 - 2.4}{6 - 2.4} = 1. \quad\checkmark
\]

\subsection{All in one cluster}
$P = \{\{1, 2, 3, 4, 5, 6\}\}$. Contingency table:
$\bigl[\begin{smallmatrix} 3 \\ 3 \end{smallmatrix}\bigr]$ (one column).
\[
  a = \binom{3}{2} + \binom{3}{2} = 6, \quad
  S_T = 6, \quad S_P = \binom{6}{2} = 15, \quad
  E[a] = \tfrac{6 \cdot 15}{15} = 6, \quad \mathrm{MaxIndex} = \tfrac{6+15}{2} = 10.5.
\]
\[
  \text{ARI} = \frac{6 - 6}{10.5 - 6} = 0. \quad\checkmark
\]
Putting everything in one cluster gets the same agreement as random ---
because it has \emph{no information} about cluster structure at all.

\subsection{Each item in its own cluster}
$P = \{\{1\}, \{2\}, \{3\}, \{4\}, \{5\}, \{6\}\}$. Contingency table:
each row has three columns with $1$ in three of the six.
\[
  a = 6 \cdot \binom{1}{2} = 0, \quad
  S_T = 6, \quad S_P = 6 \cdot \binom{1}{2} = 0, \quad
  E[a] = \tfrac{6 \cdot 0}{15} = 0, \quad \mathrm{MaxIndex} = 3.
\]
\[
  \text{ARI} = \frac{0 - 0}{3 - 0} = 0. \quad\checkmark
\]
Also gets ARI $= 0$. Singletons-everywhere is exactly as informative as
all-in-one --- you can extract no structure from either.

\subsection{A reasonable mistake}
$P = \{\{1, 2\}, \{3, 4, 5, 6\}\}$ --- we picked up the first true cluster
mostly right but lost item 3 into the merged blob with $T_2$.
\[
\text{Contingency:} \quad
\begin{array}{c | c c}
 & P_1 = \{1,2\} & P_2 = \{3,4,5,6\} \\
\hline
T_1 = \{1,2,3\} & 2 & 1 \\
T_2 = \{4,5,6\} & 0 & 3 \\
\end{array}
\]
\[
  a = \binom{2}{2} + \binom{1}{2} + \binom{0}{2} + \binom{3}{2} = 1 + 0 + 0 + 3 = 4.
\]
\[
  S_T = 6, \quad S_P = \binom{2}{2} + \binom{4}{2} = 1 + 6 = 7.
\]
\[
  E[a] = \tfrac{6 \cdot 7}{15} = 2.8, \quad \mathrm{MaxIndex} = \tfrac{6+7}{2} = 6.5.
\]
\[
  \text{ARI} = \frac{4 - 2.8}{6.5 - 2.8} = \frac{1.2}{3.7} \approx 0.324.
\]
A clearly-better-than-random but imperfect prediction. The intermediate
ARI value reflects that. Note also that RI for this case is
$\,(4 + 4)/15 \approx 0.53\,$ --- only modestly above the random baseline.
ARI translates that modest excess into a more interpretable number.

% =============================================================================
\section{Back to your K-sweep}

Recall the question that started this document: why is ARI the right
metric for the K-sweep, instead of RI?

Section~\ref{sec:problem2} answered the conceptual half. Now we can give
the quantitative answer: \textbf{ARI of a random prediction is $0$ at
every $K$}, by construction. The numerator $a - E[a]$ has expectation
zero under the null, so the ratio does too. The chance baseline that
ballooned RI from $0.5$ to $0.85$ across the K-sweep --- entirely
parasitic, entirely uninformative --- gets subtracted out completely by
ARI.

This means: when we plot ARI vs $K$ for our K-means runs, every point on
the plot is measured against the same zero. A drop from ARI $= 0.65$ at
$K=2$ to ARI $= 0.31$ at $K=12$ tells us the same story we'd see on a
fair scale: clustering quality is genuinely degrading as the author pool
grows. RI would have shown the opposite.

To make this very concrete, here are the K-sweep results from our
LessWrong 35 Authors experiment, alongside the approximate RI a random
prediction would score at each $K$ (from Table~\ref{tab:rkillertable}).

\begin{table}[H]
\centering
\small
\begin{tabular}{r r r r}
\toprule
$K$ & ARI (K-means) & random's RI & approx K-means RI \\
\midrule
2  & 0.65 & 0.500 & $\sim 0.83$ \\
4  & 0.54 & 0.631 & $\sim 0.83$ \\
6  & 0.47 & 0.729 & $\sim 0.86$ \\
8  & 0.38 & 0.787 & $\sim 0.87$ \\
10 & 0.36 & 0.825 & $\sim 0.89$ \\
12 & 0.31 & 0.852 & $\sim 0.90$ \\
\bottomrule
\end{tabular}
\caption{ARI tells the truth (clustering gets harder); RI tells the
opposite (looks like it gets easier), purely because the random baseline
grows. The K-means RI is approximate because it depends on actual cluster
sizes in each trial, but the qualitative trend is robust.}
\end{table}

If we had naively reported RI for this experiment, we would have written
``RI improves from $0.83$ to $0.90$ as we add more authors'' --- a
completely backwards conclusion that would have invalidated half of the
report.

% =============================================================================
\section{What to remember}

\begin{enumerate}\setlength\itemsep{2pt}
  \item Comparing two partitions = counting how many \emph{pairs of items}
        the partitions agree on (same-or-different).
  \item Pairs fall into four buckets: $a, b, c, d$. RI = $(a + d) /
        \binom{N}{2}$.
  \item A random clustering scores RI that is \emph{not} zero, and
        \emph{not constant} across different cluster counts $K$. The
        baseline rises with $K$ because the universe of pairs becomes
        dominated by ``trivially different''.
  \item ARI subtracts the random baseline. Conceptually:
        ARI = (RI $-$ $E[\text{RI}]$) / (1 $-$ $E[\text{RI}]$).
        Practically, the Hubert--Arabie form is
        \[
          \text{ARI} \;=\; \frac{a - E[a]}{\mathrm{MaxIndex} - E[a]},
          \qquad
          E[a] = \frac{S_T S_P}{\binom{N}{2}},
          \qquad
          \mathrm{MaxIndex} = \frac{S_T + S_P}{2}.
        \]
  \item By construction, ARI $= 0$ for random, ARI $= 1$ for perfect,
        ARI $<0$ for worse than random. These values mean the same thing
        across different $K$, different cluster sizes, different
        datasets --- which is the whole reason we use it.
\end{enumerate}

\end{document}
